\uuid{Ntjc}
\titre{Ensembles}
\auteur{bodin ; croizat ; sacre}
\organisation{exo7}
\licence{CC BY-NC-SA 4.0}
\datecreate{2022-03-08}
\niveau{L1}
\module{Algèbre}
\chapitre{Logique, ensemble, raisonnement}
\sousChapitre{Ensemble}
\difficulte{2}
\theme{Ensembles}
\exo7qcmid{hilisit.logique.033}

\contenu{
\begin{question}
\qclassification{100.02}
Soient les fonctions réelle définies par $f(x) = 3x+2$ et $g(x) = ax + b$. Pour quelle(s) valeur(s) des réels $a$ et $b$ a-t-on $g \circ f (x) = 6x + 7$ ?
\begin{answers}
	\bad{$a=1$ et $b=3$}
	\bad{$a=1$ et $b=5$}
	\good{$a=2$ et $b=3$}
	\bad{$a=2$ et $b=5$}
\end{answers}
\begin{explanations}
On calcule que $g \circ f (x) = g(3x+2) = a(3x+2) + b = 3a x + 2a + b$. Pour que cette expression soit égale à $6x+7$, on doit avoir $3a = 6$ et $2a+b=7$. Cela donne $a=2$ et $b=7- 2 \times 2 = 3$.
\end{explanations}
\end{question}
}
